\SourceLeaf{361}

\Needspace{13\baselineskip}
\begin{center}
[IV.]
\end{center}

\begin{center}
\bfseries STEREOGRAPHISCHE PROJECTION DER KUGELFLÄCHE.
\end{center}

Conjugirte complexe Zahlen sollen hier immer durch kleine und grosse gleichnamige Buchstaben bezeichnet werden.

Halbmesser der Kugel wird $=1$ gesetzt.

$n$: lineares Vergrösserungsverhältniss.

1. In dem durch die complexe Zahl $a$ bezeichneten Punkte ist
\[
n = 1+aA.
\]

2. Der gegenüberliegende Punkt wird durch die complexe Zahl $-\dfrac{1}{A}$ bezeichnet; in demselben ist also
\[
n = 1+\frac{1}{aA} = \frac{1+aA}{aA}.
\]

3. Der bicomplexen Zahl $a$ entspricht die tricomplexe
\[
\frac{2a+(1-aA)\eta}{1+aA}.
\]

4. Allgemeine Verwandlungsformel:
\[
t' = \frac{at-b}{ct+d},
\]

wo $t$, $t'$ sich auf Einerlei Punkt in der ersten oder zweiten Darstellungsform beziehen.

Hier wird
\[
n' = \text{Mod.}\ \frac{ad+bc}{(ct+d)^2}(1+tT).
\]

\Needspace{11\baselineskip}Der Minimumwerth findet statt für
\[
t = \frac{C}{D}
\]

\noindent und ist
\[
\text{Mod.}\ \frac{ad+bc}{cC+dD}\cdot\frac{D}{d} = \text{Mod.}\ \frac{ad+bc}{cC+dD}.
\]
